%=============================================================================
% NỘI DUNG PHIẾU KIỂM TRA MÔN TOÁN LỚP 9
% Thư mục: 1.SACH/LOP9/SOAN/55DE
% Lớp Toán Thầy Lĩnh - loptoanthaylinh.vn - Hotline: 0866976878
%=============================================================================

\noindent\textbf{\large I. PHẦN TRẮC NGHIỆM KHÁCH QUAN (4,0 điểm)}

\textit{Thí sinh chọn một phương án trả lời đúng nhất cho mỗi câu hỏi từ Câu 1 đến Câu 8 và ghi vào Bảng đáp án.}

\setcounterpageref{ex}{0}

%------------------------------------------------------------------------------
% CÂU 1
%------------------------------------------------------------------------------
\begin{ex}
Biểu thức $\sqrt{2x - 6}$ xác định (có nghĩa) khi và chỉ khi
\choice
{$x > 3$}
{$x \le 3$}
{\True $x \ge 3$}
{$x < 3$}
\loigiai{
Biểu thức $\sqrt{2x - 6}$ xác định khi:
$$2x - 6 \ge 0 \iff 2x \ge 6 \iff x \ge 3.$$
\textbf{Chọn C.}
}
\end{ex}

%------------------------------------------------------------------------------
% CÂU 2
%------------------------------------------------------------------------------
\begin{ex}
Giá trị của biểu thức $A = \sqrt{20} - \sqrt{45} + 3\sqrt{5}$ bằng
\choice
{$\sqrt{5}$}
{\True $2\sqrt{5}$}
{$3\sqrt{5}$}
{$4\sqrt{5}$}
\loigiai{
Ta có:
$$A = \sqrt{4 \cdot 5} - \sqrt{9 \cdot 5} + 3\sqrt{5} = 2\sqrt{5} - 3\sqrt{5} + 3\sqrt{5} = 2\sqrt{5}.$$
\textbf{Chọn B.}
}
\end{ex}

%------------------------------------------------------------------------------
% CÂU 3
%------------------------------------------------------------------------------
\begin{ex}
Trong các phương trình sau, phương trình nào là phương trình bậc nhất hai ẩn?
\choice
{$2x^2 + 3y = 5$}
{\True $3x - 4y = 7$}
{$xy + 2x = 1$}
{$0x + 0y = 6$}
\loigiai{
Phương trình bậc nhất hai ẩn $x$ và $y$ có dạng tổng quát $ax + by = c$, trong đó $a$ và $b$ không đồng thời bằng $0$.\\
Do đó, $3x - 4y = 7$ là phương trình bậc nhất hai ẩn.\\
\textbf{Chọn B.}
}
\end{ex}

%------------------------------------------------------------------------------
% CÂU 4
%------------------------------------------------------------------------------
\begin{ex}
Cặp số $(x_0; y_0)$ nào sau đây là nghiệm của hệ phương trình $\begin{cases} 2x + y = 5 \\ x - y = 1 \end{cases}$?
\choice
{$(1; 2)$}
{$(3; -1)$}
{\True $(2; 1)$}
{$(2; -1)$}
\loigiai{
Cộng từng vế hai phương trình trong hệ ta được:
$$(2x + y) + (x - y) = 5 + 1 \iff 3x = 6 \iff x = 2.$$
Thay $x = 2$ vào phương trình $x - y = 1$:
$$2 - y = 1 \iff y = 1.$$
Vậy cặp số nghiệm của hệ là $(2; 1)$.\\
\textbf{Chọn C.}
}
\end{ex}

%------------------------------------------------------------------------------
% CÂU 5
%------------------------------------------------------------------------------
\begin{ex}
Hàm số bậc nhất $y = (m - 2)x + 3$ đồng biến trên $\mathbb{R}$ khi và chỉ khi
\choice
{$m < 2$}
{$m \ge 2$}
{$m \le 2$}
{\True $m > 2$}
\loigiai{
Hàm số bậc nhất $y = ax + b$ ($a \ne 0$) đồng biến trên $\mathbb{R}$ khi $a > 0$.\\
Ở đây hệ số $a = m - 2$, do đó hàm số đồng biến khi:
$$m - 2 > 0 \iff m > 2.$$
\textbf{Chọn D.}
}
\end{ex}

%------------------------------------------------------------------------------
% CÂU 6
%------------------------------------------------------------------------------
\begin{ex}
Cho tam giác $ABC$ vuông tại $A$ có độ dài hai cạnh góc vuông $AB = 6\text{ cm}$, $AC = 8\text{ cm}$. Giá trị của $\sin B$ bằng
\choice
{$\dfrac{3}{5}$}
{\True $\dfrac{4}{5}$}
{$\dfrac{3}{4}$}
{$\dfrac{4}{3}$}
\loigiai{
Áp dụng định lí Pythagore trong tam giác vuông $ABC$:
$$BC = \sqrt{AB^2 + AC^2} = \sqrt{6^2 + 8^2} = \sqrt{100} = 10\text{ cm}.$$
Khi đó:
$$\sin B = \dfrac{\text{cạnh đối}}{\text{cạnh huyền}} = \dfrac{AC}{BC} = \dfrac{8}{10} = \dfrac{4}{5}.$$
\textbf{Chọn B.}
}
\end{ex}

%------------------------------------------------------------------------------
% CÂU 7
%------------------------------------------------------------------------------
\begin{ex}
Cho đường tròn $(O)$ có góc nội tiếp $\widehat{BAC} = 50^\circ$ chắn cung nhỏ $BC$. Số đo góc ở tâm $\widehat{BOC}$ chắn cung $BC$ là
\choice
{$50^\circ$}
{$25^\circ$}
{\True $100^\circ$}
{$130^\circ$}
\loigiai{
Trong một đường tròn, số đo của góc nội tiếp bằng nửa số đo góc ở tâm cùng chắn một cung:
$$\widehat{BAC} = \dfrac{1}{2} \widehat{BOC} \implies \widehat{BOC} = 2 \cdot \widehat{BAC} = 2 \cdot 50^\circ = 100^\circ.$$
\textbf{Chọn C.}
}
\end{ex}

%------------------------------------------------------------------------------
% CÂU 8
%------------------------------------------------------------------------------
\begin{ex}
Gieo một con xúc xắc cân đối và đồng chất $6$ mặt một lần. Xác suất để xuất hiện mặt có số chấm là số nguyên tố bằng
\choice
{$\dfrac{1}{6}$}
{$\dfrac{1}{3}$}
{\True $\dfrac{1}{2}$}
{$\dfrac{2}{3}$}
\loigiai{
Tập hợp các kết quả có thể xảy ra khi gieo xúc xắc là:
$$\Omega = \{1; 2; 3; 4; 5; 6\} \implies n(\Omega) = 6.$$
Các mặt có số chấm là số nguyên tố là:
$$A = \{2; 3; 5\} \implies n(A) = 3.$$
Xác suất của biến cố là:
$$P(A) = \dfrac{n(A)}{n(\Omega)} = \dfrac{3}{6} = \dfrac{1}{2}.$$
\textbf{Chọn C.}
}
\end{ex}

\vspace{0.3cm}
\noindent\textbf{\large II. PHẦN TỰ LUẬN (6,0 điểm)}

%------------------------------------------------------------------------------
% BÀI 1 (TỰ LUẬN)
%------------------------------------------------------------------------------
\begin{ex}[1,5 điểm]
Giải hệ phương trình sau bằng \textbf{phương pháp cộng đại số}:
$$\begin{cases} 3x - 2y = 7 \quad &(1) \\ 2x + 2y = 8 \quad &(2) \end{cases}$$
\loigiai{
Nhận thấy hệ số của ẩn $y$ trong hai phương trình đối nhau ($-2$ và $2$), ta cộng từng vế của phương trình (1) và phương trình (2):
\begin{align*}
(3x - 2y) + (2x + 2y) &= 7 + 8 \\
5x &= 15 \\
x &= 3.
\end{align*}
Thay $x = 3$ vào phương trình (2):
\begin{align*}
2 \cdot 3 + 2y &= 8 \\
6 + 2y &= 8 \\
2y &= 8 - 6 \\
2y &= 2 \\
y &= 1.
\end{align*}
Vậy hệ phương trình đã cho có nghiệm duy nhất là $(x; y) = (3; 1)$.
}
\end{ex}

%------------------------------------------------------------------------------
% BÀI 2 (TỰ LUẬN)
%------------------------------------------------------------------------------
\begin{ex}[1,5 điểm]
Giải hệ phương trình sau bằng \textbf{phương pháp cộng đại số}:
$$\begin{cases} 3x + 5y = 1 \quad &(1) \\ 2x - 3y = 7 \quad &(2) \end{cases}$$
\loigiai{
Nhân hai vế của phương trình (1) với $3$ và nhân hai vế của phương trình (2) với $5$, ta được hệ phương trình tương đương:
$$\begin{cases} 9x + 15y = 3 \quad &(3) \\ 10x - 15y = 35 \quad &(4) \end{cases}$$
Cộng từng vế của phương trình (3) và phương trình (4), ta được:
\begin{align*}
(9x + 15y) + (10x - 15y) &= 3 + 35 \\
19x &= 38 \\
x &= 2.
\end{align*}
Thay $x = 2$ vào phương trình (1):
\begin{align*}
3 \cdot 2 + 5y &= 1 \\
6 + 5y &= 1 \\
5y &= 1 - 6 \\
5y &= -5 \\
y &= -1.
\end{align*}
Vậy hệ phương trình đã cho có nghiệm duy nhất là $(x; y) = (2; -1)$.
}
\end{ex}

%------------------------------------------------------------------------------
% BÀI 3 (TỰ LUẬN)
%------------------------------------------------------------------------------
\begin{ex}[1,5 điểm]
Cho biểu thức:
$$P = \left( \dfrac{\sqrt{x}}{\sqrt{x}-1} - \dfrac{1}{x-\sqrt{x}} \right) : \left( \dfrac{1}{\sqrt{x}+1} + \dfrac{2}{x-1} \right) \quad \text{với } x > 0 \text{ và } x \ne 1.$$
\begin{enumerate}[a)]
	\item Rút gọn biểu thức $P$.
	\item Tính giá trị của biểu thức $P$ khi $x = 4$.
\end{enumerate}
\loigiai{
Với điều kiện $x > 0$ và $x \ne 1$, ta có:
\begin{enumerate}[a)]
	\item Rút gọn biểu thức $P$:
	
	Xét biểu thức trong ngoặc thứ nhất:
	\begin{align*}
	A &= \dfrac{\sqrt{x}}{\sqrt{x}-1} - \dfrac{1}{\sqrt{x}(\sqrt{x}-1)} \\
	&= \dfrac{\sqrt{x} \cdot \sqrt{x} - 1}{\sqrt{x}(\sqrt{x}-1)} = \dfrac{x - 1}{\sqrt{x}(\sqrt{x}-1)} \\
	&= \dfrac{(\sqrt{x}-1)(\sqrt{x}+1)}{\sqrt{x}(\sqrt{x}-1)} = \dfrac{\sqrt{x}+1}{\sqrt{x}}.
	\end{align*}
	
	Xét biểu thức trong ngoặc thứ hai:
	\begin{align*}
	B &= \dfrac{1}{\sqrt{x}+1} + \dfrac{2}{(\sqrt{x}-1)(\sqrt{x}+1)} \\
	&= \dfrac{\sqrt{x}-1+2}{(\sqrt{x}-1)(\sqrt{x}+1)} = \dfrac{\sqrt{x}+1}{(\sqrt{x}-1)(\sqrt{x}+1)} \\
	&= \dfrac{1}{\sqrt{x}-1}.
	\end{align*}
	
	Do đó:
	\begin{align*}
	P &= A : B = \dfrac{\sqrt{x}+1}{\sqrt{x}} : \dfrac{1}{\sqrt{x}-1} \\
	&= \dfrac{\sqrt{x}+1}{\sqrt{x}} \cdot (\sqrt{x}-1) = \dfrac{(\sqrt{x}+1)(\sqrt{x}-1)}{\sqrt{x}} \\
	&= \dfrac{x-1}{\sqrt{x}}.
	\end{align*}
	Vậy $P = \dfrac{x-1}{\sqrt{x}}$.
	
	\item Tính giá trị của $P$ khi $x = 4$:
	
	Thay $x = 4$ (thỏa mãn điều kiện $x > 0$ và $x \ne 1$) vào biểu thức rút gọn $P$:
	$$P = \dfrac{4 - 1}{\sqrt{4}} = \dfrac{3}{2}.$$
	Vậy khi $x = 4$ thì $P = \dfrac{3}{2}$.
\end{enumerate}
}
\end{ex}

%------------------------------------------------------------------------------
% BÀI 4 (TỰ LUẬN)
%------------------------------------------------------------------------------
\begin{ex}[1,5 điểm]
\textbf{(Giải bài toán bằng cách lập hệ phương trình)}

Hai tổ công nhân cùng làm chung một đoạn đường thì sau $12$ ngày sẽ hoàn thành. Nếu tổ thứ nhất làm riêng trong $4$ ngày, sau đó tổ thứ hai làm tiếp một mình trong $10$ ngày thì cả hai tổ hoàn thành được $\dfrac{1}{2}$ công việc. Hỏi nếu làm riêng thì mỗi tổ hoàn thành đoạn đường đó trong bao nhiêu ngày?
\loigiai{
Gọi thời gian tổ thứ nhất làm một mình hoàn thành toàn bộ đoạn đường là $x$ (ngày).\\
Gọi thời gian tổ thứ hai làm một mình hoàn thành toàn bộ đoạn đường là $y$ (ngày).\\
Điều kiện: $x > 12, y > 12$.

Trong $1$ ngày:
\begin{itemize}
	\item Tổ thứ nhất làm được: $\dfrac{1}{x}$ (đoạn đường).
	\item Tổ thứ hai làm được: $\dfrac{1}{y}$ (đoạn đường).
\end{itemize}

Vì cả hai tổ cùng làm thì hoàn thành đoạn đường trong $12$ ngày nên trong $1$ ngày cả hai tổ làm được $\dfrac{1}{12}$ đoạn đường. Ta có phương trình:
$$\dfrac{1}{x} + \dfrac{1}{y} = \dfrac{1}{12} \quad (1)$$

Nếu tổ thứ nhất làm riêng trong $4$ ngày, sau đó tổ thứ hai làm riêng trong $10$ ngày thì hoàn thành được $\dfrac{1}{2}$ công việc, nên ta có phương trình:
$$4 \cdot \dfrac{1}{x} + 10 \cdot \dfrac{1}{y} = \dfrac{1}{2} \quad (2)$$

Từ (1) và (2), ta có hệ phương trình:
$$\begin{cases} \dfrac{1}{x} + \dfrac{1}{y} = \dfrac{1}{12} \\[8pt] \dfrac{4}{x} + \dfrac{10}{y} = \dfrac{1}{2} \end{cases}$$

Đặt $u = \dfrac{1}{x}$ và $v = \dfrac{1}{y}$ ($u > 0, v > 0$), hệ phương trình trở thành:
$$\begin{cases} u + v = \dfrac{1}{12} \\[6pt] 4u + 10v = \dfrac{1}{2} \end{cases}$$

Nhân hai vế của phương trình thứ nhất với $4$:
$$\begin{cases} 4u + 4v = \dfrac{1}{3} \\[6pt] 4u + 10v = \dfrac{1}{2} \end{cases}$$

Trừ từng vế của phương trình thứ hai cho phương trình thứ nhất:
$$6v = \dfrac{1}{2} - \dfrac{1}{3} \iff 6v = \dfrac{1}{6} \iff v = \dfrac{1}{36}.$$

Thay $v = \dfrac{1}{36}$ vào phương trình $u + v = \dfrac{1}{12}$:
$$u = \dfrac{1}{12} - \dfrac{1}{36} = \dfrac{3}{36} - \dfrac{1}{36} = \dfrac{2}{36} = \dfrac{1}{18}.$$

Suy ra:
$$\begin{cases} \dfrac{1}{x} = \dfrac{1}{18} \implies x = 18 \quad \text{(thỏa mãn điều kiện)} \\[8pt] \dfrac{1}{y} = \dfrac{1}{36} \implies y = 36 \quad \text{(thỏa mãn điều kiện)} \end{cases}$$

Vậy nếu làm riêng:
\begin{itemize}
	\item Tổ thứ nhất hoàn thành đoạn đường trong $18$ ngày.
	\item Tổ thứ hai hoàn thành đoạn đường trong $36$ ngày.
\end{itemize}
}
\end{ex}
